\documentclass[a4paper,11pt]{article}
\usepackage[utf8]{inputenc}
\usepackage[T1]{fontenc}
\usepackage[spanish]{babel}
\usepackage{geometry}
\usepackage{array}
\usepackage{colortbl}
\usepackage{xcolor}
\usepackage{booktabs}
\usepackage{enumitem}
\usepackage{fancyhdr}
\usepackage{tikz}

\geometry{left=2cm, right=2cm, top=2.5cm, bottom=2cm}

\definecolor{violeta}{HTML}{7B2D8E}
\definecolor{azulg}{HTML}{3498DB}
\definecolor{rojog}{HTML}{E74C3C}
\definecolor{gris}{HTML}{F2F2F2}
\definecolor{verde}{HTML}{27AE60}
\definecolor{naranja}{HTML}{E67E22}

\pagestyle{fancy}
\fancyhf{}
\fancyhead[L]{\small\textbf{Nuevo Compañeros 1 — U03: La Familia}}
\fancyhead[R]{\small\textbf{Píldora formativa 3.8}}
\fancyfoot[C]{\thepage}

\setlength{\parindent}{0pt}

\begin{document}

\begin{center}
{\LARGE\textbf{PÍLDORA FORMATIVA 3.8}}\\[0.3cm]
{\Large Comprensión y producción escrita: el correo de Marta}\\[0.2cm]
{\normalsize Sección Destrezas (p.40) — Proyectable — 3 diapositivas}\\[0.2cm]
{\small Versión enriquecida secuencial para Fases L4--E5}
\end{center}

\vspace{0.5cm}

\noindent\fbox{\parbox{\dimexpr\linewidth-2\fboxsep-2\fboxrule}{%
\textbf{Contenido para el profesor}\\[0.2cm]
El correo de Marta a Pierre (p.40) es el texto central de la sección de Destrezas. Integra todos los contenidos de la unidad: vocabulario de parentesco (padre, madre, hermano Mario 11 años, hermana Ana 6 años), verbos regulares (trabajar, estudiar, vivir, salir), verbo tener, posesivos y funciones comunicativas.\\[0.2cm]
El correo tiene una estructura clara que el estudiante replicará en la actividad 3: saludo, párrafo familia, párrafo amigos, párrafo instituto/asignaturas, despedida. Esta píldora marca la transición de leer a escribir: el texto leído se convierte en el worked example para la producción escrita.
}}

\vspace{0.8cm}

%% ============================================================
%% DIAPOSITIVA 1
%% ============================================================

\noindent\textcolor{violeta}{\rule{\linewidth}{2.5pt}}

\section*{DIAPOSITIVA 1 — ANALIZA EL MODELO: ESTRUCTURA Y CONTENIDO}

{\small Fase L4 (Post-lectura) + E1 (Modelo / Worked example) — Técnica: Texto cromático -- análisis de estructura (azul) vs. contenido personal (rojo)}\\
{\small\textit{Principio:} El estudiante ya comprende el correo; ahora lo analiza como modelo para su escritura. Con dos colores se separa lo que debe mantener (estructura) de lo que debe cambiar (datos personales). Esto reduce la carga cognitiva: el estudiante no parte de cero.}

\vspace{0.5cm}

\subsection*{En pantalla}

\begin{center}
{\Large\textbf{ANALIZA EL CORREO DE MARTA ANTES DE ESCRIBIR EL TUYO}}
\end{center}

\vspace{0.3cm}

El correo de Marta con marcas de color:

\vspace{0.3cm}

\noindent\textbf{\textcolor{azulg}{AZUL (estructura = mantener):}}
\begin{itemize}[leftmargin=2cm, itemsep=2pt]
\item \textcolor{azulg}{Hola, [nombre]! ¿Qué tal estás?}
\item \textcolor{azulg}{Hoy te hablo de mi familia:}
\item \textcolor{azulg}{Y en el instituto, ¿qué tal el nuevo curso?}
\item \textcolor{azulg}{Un saludo desde [ciudad]!}
\end{itemize}

\vspace{0.2cm}

\noindent\textbf{\textcolor{rojog}{ROJO (contenido personal = cambiar):}}
\begin{itemize}[leftmargin=2cm, itemsep=2pt]
\item \textcolor{rojog}{mi padre trabaja en un hotel} $\rightarrow$ ¿Dónde trabajan TUS padres?
\item \textcolor{rojog}{mi hermano Mario tiene once años} $\rightarrow$ ¿Tienes hermanos? ¿Cómo se llaman?
\item \textcolor{rojog}{tengo una compañera nueva: se llama Bárbara} $\rightarrow$ ¿Quiénes son tus amigos?
\item \textcolor{rojog}{tenemos muchos deberes de Tecnología} $\rightarrow$ ¿Qué asignaturas tienes?
\item \textcolor{rojog}{A mí me gusta mucho la Historia} $\rightarrow$ ¿Cuál es TU asignatura favorita?
\end{itemize}

\vspace{0.3cm}

Esquema visual: 5 bloques del correo con flechas que indican qué cambiar.

\vspace{0.5cm}

\subsection*{Instrucciones para el profesor}

\begin{enumerate}[leftmargin=1.5cm, itemsep=4pt]
\item Con la herramienta marcador de BlinkLearning (o en la pizarra), marque el correo con dos colores: AZUL = estructura (mantener) / ROJO = contenido (cambiar con vuestros datos).
\item Lea el saludo: \textit{Hola, Pierre! ¿Qué tal estás? ¿Esto lo mantenéis o lo cambiáis?} Lo mantenéis, pero con el nombre de vuestro amigo. Es azul.
\item Lea: \textit{Mi padre trabaja en un hotel. ¿Esto?} Lo cambiáis con la profesión de vuestro padre o vuestra madre. Es rojo.
\item Continúe párrafo a párrafo. Pida que los propios estudiantes decidan azul o rojo para cada frase.
\item Reparta la tarjeta de estrategia — Escribir un correo (Caja 3, tarjeta B). Señale los 5 pasos y el Truco del semáforo: ROJO = planifica, AMARILLO = escribe, VERDE = revisa.
\end{enumerate}

\newpage

%% ============================================================
%% DIAPOSITIVA 2
%% ============================================================

\noindent\textcolor{violeta}{\rule{\linewidth}{2.5pt}}

\section*{DIAPOSITIVA 2 — ESCRIBE TU CORREO: PASO A PASO}

{\small Fase E2--E3 (Planificación + Borrador) — Técnica: Writing frame con pasos del libro + producción guiada con apoyo visible}\\
{\small\textit{Principio:} El estudiante sigue los 6 pasos del libro para escribir su correo. No parte de cero: tiene el modelo analizado, las tarjetas de vocabulario, los esquemas comunicativos y las tarjetas de estrategia.}

\vspace{0.5cm}

\subsection*{En pantalla}

Los 6 pasos del libro:

\begin{enumerate}[leftmargin=1.5cm, itemsep=4pt]
\item Planifica tu correo: decide de qué vas a hablar en cada párrafo. (Mira el correo de Marta.)
\item Escribe una corta descripción de tu familia.
\item Escribe sobre los compañeros y amigos de tu clase.
\item Escribe sobre las asignaturas y el horario de este curso.
\item No te olvides de escribir: un saludo para empezar y una frase de despedida.
\item Repasa y comprueba las descripciones de tu familia y amigos y la forma correcta de los verbos.
\end{enumerate}

\vspace{0.3cm}

Writing frame con inicio de cada párrafo:

\begin{center}
\renewcommand{\arraystretch}{2}
\begin{tabular}{|>{\raggedright\arraybackslash}p{13cm}|}
\hline
\rowcolor{azulg!10}
Hola, \rule{4cm}{0.4pt}! ¿Qué tal estás? \\
\hline
Hoy te hablo de mi familia: mi padre/madre \rule{6cm}{0.4pt} \\
\hline
En mi clase tengo muchos amigos: \rule{6cm}{0.4pt} \\
\hline
Este curso \rule{8cm}{0.4pt} \\
\hline
\rowcolor{azulg!10}
Un saludo/abrazo desde \rule{4cm}{0.4pt}! \\
\hline
\end{tabular}
\end{center}

\vspace{0.5cm}

\subsection*{Instrucciones para el profesor}

\begin{enumerate}[leftmargin=1.5cm, itemsep=4pt]
\item Lea los 6 pasos del libro en voz alta.
\item Señale el writing frame: \textit{Aquí tenéis el inicio de cada párrafo. Completadlo con vuestros datos.}
\item Recuerde: \textit{Tenéis todas las tarjetas disponibles. Si necesitáis vocabulario de familia, mirad la Caja 1. Si necesitáis conjugar un verbo, mirad la Caja 2.}
\item Dé 5--7 minutos para el borrador. Circule y atienda dudas individualmente. Si un estudiante está bloqueado, señale el modelo: \textit{Mira qué escribe Marta aquí. ¿Qué puedes escribir tú?}
\item Si la actividad se completa en casa, los pasos 1--3 del libro se trabajan en clase y los pasos 4--6 se asignan como tarea.
\end{enumerate}

\newpage

%% ============================================================
%% DIAPOSITIVA 3
%% ============================================================

\noindent\textcolor{violeta}{\rule{\linewidth}{2.5pt}}

\section*{DIAPOSITIVA 3 — REVISA CON TU COMPAÑERO}

{\small Fase E4--E5 (Revisión entre pares + Versión final) — Técnica: Peer response checklist + focused corrective feedback}\\
{\small\textit{Principio:} El compañero lee el borrador y marca con una checklist sencilla (sí/no). Solo se revisa un tipo de error por ronda para no sobrecargar. El estudiante corrige y produce la versión final.}

\vspace{0.5cm}

\subsection*{En pantalla}

Checklist de revisión (el compañero marca sí o no):

\begin{center}
\renewcommand{\arraystretch}{1.6}
\begin{tabular}{|>{\raggedright\arraybackslash}p{10cm}|>{\centering\arraybackslash}p{1.5cm}|>{\centering\arraybackslash}p{1.5cm}|}
\hline
\rowcolor{verde!15}
\textbf{Pregunta} & \textbf{Sí} & \textbf{No} \\
\hline
1. ¿Tiene saludo (Hola, [nombre])? & & \\
\hline
2. ¿Habla de la familia (padres, hermanos)? & & \\
\hline
3. ¿Habla de amigos o compañeros? & & \\
\hline
4. ¿Habla de asignaturas o el instituto? & & \\
\hline
5. ¿Tiene despedida (Un saludo / Un abrazo)? & & \\
\hline
6. ¿Los verbos están bien escritos (trabaja, tiene, estudia)? & & \\
\hline
\end{tabular}
\end{center}

\vspace{0.3cm}

Instrucción: \textit{Intercambia tu correo con tu compañero. Lee su correo y marca la checklist. Si algo falta, díselo.}

\vspace{0.3cm}

Tras la revisión: \textit{Corrige tu correo y escribe la versión final.}

\vspace{0.5cm}

\subsection*{Instrucciones para el profesor}

\begin{enumerate}[leftmargin=1.5cm, itemsep=4pt]
\item Forme parejas (diferentes a las de lectura si es posible).
\item Explique la checklist: \textit{Vais a leer el correo de vuestro compañero. No es un examen: es una ayuda. Marcad sí o no en cada punto.}
\item Dé 3 minutos para la lectura + checklist.
\item Los estudiantes devuelven el correo con la checklist. Cada uno corrige su borrador (2--3 minutos).
\item Opcional: si hay tiempo, pida a 2--3 voluntarios que lean su correo en voz alta. Use reformulaciones correctivas para los errores gramaticales.
\item Si la versión final se completa en casa, indique: \textit{Corregid el correo en casa siguiendo la checklist de vuestro compañero. La próxima clase, 2--3 voluntarios leerán su correo.}
\end{enumerate}

\end{document}
